\chapter{Ergon Informatica S.r.l}
\label{chap:introduzione}


\textit{Il presente capitolo delinea il contesto operativo all'interno del quale si è svolta l'attività di tirocinio. Attraverso un'analisi della struttura aziendale e delle metodologie di sviluppo adottate, si intende definire il quadro di riferimento per comprendere le successive decisioni progettuali, il contesto in cui opera e l'importanza verso l'innovazione e il miglioramento continuo. Le informazioni qui riportate sono state raccolte mediante l'osservazione attiva dei flussi di lavoro, l'analisi della documentazione tecnica interna e il confronto con le figure di riferimento aziendale.}

\section{Profilo aziendale}
In questa sezione viene descritta l'identità di Ergon Informatica S.r.l., ripercorrendone brevemente la storia e il posizionamento strategico nel panorama IT. Sono analizzati i principali domini applicativi di riferimento e la tipologia di clientela servita, fornendo una visione d'insieme atta a contestualizzare la rilevanza del tirocinio.


\section{Metodologie di sviluppo e processi operativi}
Viene analizzato il ciclo di vita del \textit{software} adottato in azienda, mettendo in luce le pratiche metodologiche che guidano il team di sviluppo dal recepimento dei requisiti fino al rilascio. Particolare attenzione è rivolta agli strumenti tecnologici utilizzati e alle procedure di manutenzione, supportate da schemi e flussi che testimoniano l'approccio rigoroso dell'azienda.

\section{Innovazione tecnologica e cultura aziendale}
La seguente sezione esplora la propensione di Ergon verso l'innovazione ed il miglioramento continuo. Attraverso l'analisi delle scelte tecnologiche e della cultura aziendale, si evidenzia come tali elementi abbiano influito direttamente sul progetto di tirocinio, ponendo le basi per l'analisi dei capitoli successivi.

\newpage